\subsection{The cuspidal fibre: the precise stacky correction}

\begin{proposition}[\status{Refuted}: roots cannot repair a meridian kernel]\label{P1e-roots}
Let $(X,D)$ be a smooth complex normal-crossings pair, and let a finite root
construction along its branches leave $U=X\setminus D$ unchanged.
If the old meridian homomorphism $\rho:\Z^r\to\pi_1(U)$ has nonzero kernel,
the corresponding meridian map of the root stack also has nonzero kernel.
More generally this holds for any finite Kummer change of the local lattice.
Finite local monodromy of a family does not by itself force failure of the
criterion on smooth models.
\end{proposition}
\begin{proof}
In a local chart for roots of orders $n_1,\ldots,n_r$, write $z_i=w_i^{n_i}$.
A normal circle in the chart maps to $n_i$ old normal circles.  Thus the new
map is $\rho\circ\operatorname{diag}(n_1,\ldots,n_r)$.  The image of this
diagonal matrix is a finite-index sublattice of $\Z^r$; it meets every nonzero
subgroup of $\Z^r$ nontrivially.  This proves the assertion, and the same
argument applies to any finite-index Kummer sublattice.  The displayed integer
kernel remains a kernel after profinite completion.  A root of the entire
nonreduced fibre may itself be singular; merely calling that operation a root
does not bypass this lattice obstruction.

For finite monodromy without failure, take $F\times\Delta$ with boundary
$F\times\{0\}$, where $F$ is an affine hyperbolic curve.  Both strata are
aspherical and the normal circle is the direct $\Z$ factor of
$\pi_1(F)\times\Z$.  Its monodromy is the identity.  For nontrivial finite
monodromy choose a free finite-order automorphism $\phi$ of a hyperbolic curve
$F$ (for example a deck transformation of an unramified cyclic cover), and
let $\mu_m$ act on $F\times\Delta_s$ by $(p,s)\mapsto(\phi(p),\zeta s)$.
The total quotient is a smooth manifold because the action on $F$ is free.
Its reduced central fibre is $F/\mu_m$, and its normal circle projects to
$m$ times the generator around $t=s^m=0$.  It therefore injects.  The open
space is a surface bundle over $\Delta^*$, and the closed stratum is a
hyperbolic curve, so the criterion holds on this smooth model.
\end{proof}

\begin{example}[\status{Proved}: all lattices and mapping spaces of the cusp]\label{P1e-cusp}
Consider $f(x,y)=y^2-x^3$ and the family $f=t$ over a sufficiently small disk,
with the affine point at infinity of each elliptic curve omitted.  It has
pointed type $(1,1)$, not the unmarked proper type $(1,0)$.
The minimal embedded normal-crossings resolution of $f^{-1}(0)$ has four
components $C,E_2,E_3,E_6$ of multiplicities $1,2,3,6$, with $E_6$ meeting
each of the other three once and no other intersections.
The open complement has
\[
 G=\pi_1(U)=\langle a,b\mid a^2=b^3\rangle\cong B_3,\qquad
 z=a^2=b^3,\qquad m=ab^{-1}.
\]
With paths chosen compatibly, the meridians of $E_3,E_2,E_6,C$ are respectively
$a,b,z,m$.  All rank-one meridian maps inject.  At the three nodes, written
in the order $(E_6,E_j)$, the rank-two maps have the following kernels:
\[
 \begin{array}{c|c|c}
 \text{node}&(r,s)\longmapsto&\text{kernel}\\ \hline
 E_6\cap E_2&z^r b^s&\Z(1,-3)\\
 E_6\cap E_3&z^r a^s&\Z(1,-2)\\
 E_6\cap C&z^r m^s&0 .
 \end{array}                                                   \tag{CU}
\]
All strata are aspherical, but the criterion fails at the first two nodes.
Taking one object per connected stratum gives eight objects: the open object
with automorphism group $G$; four component objects with groups
$\Z,1,1,F_2$ for $C,E_2,E_3,E_6$; and three node objects with trivial groups.
For a component object, its mapping space to the open object is the discrete
coset set $G/\langle m_j\rangle$, with $m_j=m,b,a,z$ respectively.
At the first two nodes the mapping spaces to the open object are disjoint
unions of circles indexed by $G/\langle b\rangle$ and $G/\langle a\rangle$.
At the last node the mapping space is the discrete set
$G/\langle z,m\rangle$.
The maps from a node to the incident $E_2$ or $E_3$ object form a circle;
the maps to $E_6$ are $F_2/\langle c_j\rangle$, where $c_j$ is the appropriate
boundary word of the three-punctured sphere; the maps from $E_6\cap C$ to
$C$ form a singleton.  All other distinct-stratum mapping spaces are empty.
The group actions and compositions are induced by the inclusions of these
normal links; choices of the displayed paths change the descriptions by
compatible conjugations.
\end{example}
\begin{proof}
The three blow-up charts relevant to the successive singular points are
\[
 f=x^2(v^2-x),\qquad
 f=u^2v^3(v-u),\qquad
 f=u^6w^3(w-1).
\]
They follow from $y=xv$, then $x=uv$, then $v=uw$.  The last chart shows the
multiplicity six and the three distinct intersection points $w=0,1,\infty$.
The other charts show the strict components of multiplicity two and three.
After the first blow-up the strict cusp is tangent to $E_2$; after the second
there is a triple point; the third is therefore necessary for a normal-crossings
divisor on a surface.  These polynomial identities are checked symbolically
in \texttt{compactifications\_checks.py}.

The weighted circle action $(x,y)\mapsto(e^{2i\theta}x,e^{3i\theta}y)$ has
regular orbit $z$ and exceptional orbit multiplicities two and three.
Cutting its torus-knot exterior along the separating annulus gives the
amalgamated presentation $\langle a,b\mid a^2=b^3\rangle$: the annulus core
is twice the core of one solid torus and three times the other.  Equivalently,
this is the complement of the $(2,3)$ torus knot, since intersecting
$x=s^2,y=s^3$ with a small sphere gives that knot.  Weighted radial contraction
identifies its complement homotopy type with $f^{-1}(\Delta^*)$.
The degree homomorphism induced by $f$ is
\[
 \deg(a)=3,\quad\deg(b)=2,\quad\deg(z)=6,\quad\deg(m)=1.
\]
In a chart near the unmodified end of $E_2$, $f=y^2(1-u^3y)$, so its meridian
is the exceptional orbit of degree two, namely $b$.  At the other end,
$f=x^3(w^2x-1)$ identifies the $E_3$ meridian with $a$.  In the last blow-up
chart, $(x,y)=(u^2w,u^3w^2)$, rotation of $u$ is exactly the weighted circle
$z$.  The meridian $m=ab^{-1}=a^{-1}b^2$ of the strict cusp is the usual
peripheral meridian; the two expressions agree because $z$ is central.
The centrality of $z$ follows directly from the amalgam presentation.

The first two kernels in (CU) now follow from $z=b^3=a^2$ and the infinite
orders certified by the degree homomorphism.  For the last one, quotient by
the central $\langle z\rangle$.  The result is $C_2*C_3$, in which the image
of $m=ab^{-1}$ is a cyclically reduced alternating word and has infinite
order.  Thus $z^rm^s=1$ forces $s=0$, then $r=0$ by the degree homomorphism.
The program independently checks the representation
\[
 a\longmapsto\begin{pmatrix}0&-1\\1&0\end{pmatrix},\qquad
 b\longmapsto\begin{pmatrix}0&-1\\1&1\end{pmatrix},
\]
where $a^2=b^3=-I$ and $m$ is nontrivial unipotent.  Combined with the degree
homomorphism this also proves the last lattice injective.

The component strata are $C^\circ\simeq\C^*$,
$E_2^\circ\simeq\C$, $E_3^\circ\simeq\C$, and
$E_6^\circ\simeq\PP^1\setminus\{0,1,\infty\}$.
The open space fibres over a circle with once-punctured torus fibre and is
aspherical.  For a homomorphism from a torus lattice to an aspherical target,
the mapping-space homotopy fibre has components $K(\ker\rho,1)$ indexed by
the cosets of its image.  Applying this to (CU), and to the rank-one links
inside the four component strata, gives every displayed mapping space.
In particular the failed adjacent maps into $\C$ have a whole circle as
mapping space, while the adjacent maps into the three-punctured sphere have
injective peripheral cyclic subgroups.  This accounts for all incidence
relations of the four-component star.
\end{proof}

\begin{theorem}[\status{Proved}: stable reduction supplies a different good stack]\label{P1e-cusp-stack}
For the cusp family, no finite roots along the resolved fibre repair the two
failed lattices in (CU).  A good stacky compactification of the same open
family is instead
\[
 \mathcal X=[(E^\circ\times\Delta_s)/\mu_6],\qquad
 E^\circ=\{Y^2=X^3+1\},
\]
with action
\[
 (X,Y,s)\longmapsto(\zeta^{-2}X,\zeta^{-3}Y,\zeta s),
 \qquad t=s^6,quad x=s^2X,quad y=s^3Y.
\]
Its boundary is the smooth stack $[E^\circ/\mu_6]$.  It is aspherical, with
fundamental group $C_2*C_3$, and its normal-circle injection is
$\Z\langle z\rangle\hookrightarrow B_3$.
The exit category is ordinary, with two stratum objects, automorphism groups
$C_2*C_3$ and $B_3$, and crossing Hom-set $B_3/\langle z\rangle$.
The two actions on this set are the left regular $B_3$ action and the right
action through $B_3/\langle z\rangle\cong C_2*C_3$, with composition fixed by
these actions.
\end{theorem}
\begin{proof}
For $s\ne0$ the coordinate transformation is invertible after choosing a sixth
root of $t$, and two choices differ by precisely the displayed action.
Consequently its quotient open substack is the original $U$.
The action is free on $s\ne0$, so no extra open stabilizers have been introduced.
The stack itself is smooth because its atlas is smooth; its boundary is a
smooth divisor in this atlas.

The quotient of the affine elliptic curve by $\mu_6$ has coarse coordinate
$w=Y^2=X^3+1$.  At $w=0$ the stabilizer is $\mu_2$, and at $w=1$ it is
$\mu_3$; there is no other nontrivial stabilizer.  Its orbifold fundamental
group is therefore $C_2*C_3$.  The finite \'etale atlas $E^\circ$ is a
once-punctured torus and aspherical; the universal orbifold cover of the
quotient is consequently contractible.  Rotation of $s$ once at a fixed
point of the atlas gives the weighted rotation $(x,y)=(s^2X,s^3Y)$, hence the
normal meridian $z$, which has infinite order by degree six.  The open stratum
is already aspherical, so the criterion applies, and its mapping spaces are
the stated discrete cosets.  The first assertion is
Proposition~\ref{P1e-roots} applied to the two nonzero kernels of (CU).

This construction combines finite stable base change, a new model, and the
quotient stack retaining the finite automorphism action.  It is not the root
stack of the resolved boundary.  In a general finite-monodromy stable
reduction of order $m$, the analogous quotient of the equivariant stable model
by $\mu_m$ is the relevant stack structure.  A smooth coarse model may already
suffice when that action is free, as shown in Proposition~\ref{P1e-roots}.
All statements here are over $\C$; the smooth tame $\mu_6$ discussion does not
extend unchanged to residue characteristics two and three.
\end{proof}
